\documentclass[12pt]{article}
\usepackage[utf8]{inputenc}
\usepackage[T1]{fontenc}
\usepackage{hyperref}
\title{Vori Language Resources for K-12 Curriculum Development}
\author{Compiled by Claude Sonnet 5.5 (claude-sonnet-5-5) for Andr\'as Kornai}
\date{}
\begin{document}
\maketitle

\section{Introduction}
Vori (ISO 639-3 sde; Glottocode suru1258), also called Surubu or Skrubu in the literature, is an East Kainji (Niger-Congo) language of Kaduna State, Nigeria, with its main centre at Geshere, west of Mariri on the Lere road. Roger Blench (2016) reported that the people called Surubu in earlier literature call themselves Vori. Speaker figures conflict: 3,000 (Wikipedia, 2016), 7,200 (survey reference), and 50,000--70,000 in another entry. Do not confuse Vori with Voro (ISO vor), an Adamawa language. Searches found no educational materials at all.

\section{Primary Educational Resources}
None found.

\section{Linguistic and Reference Materials}
\begin{itemize}
\item Wikipedia article on Vori: \url{https://en.wikipedia.org/wiki/Vori_language}
\item Glottolog entry (status ``shifting''): \url{https://glottolog.org/resource/languoid/id/suru1258}
\item ASJP Surubu wordlist, compiled from Williamson's comparative lists: \url{https://asjp.clld.org/languages/SURUBU}
\item Wikipedia on the Kainji languages: \url{https://en.wikipedia.org/wiki/Kainji_languages}
\item Wikipedia on East Kainji: \url{https://www.wikipedia.com/wiki/East_Kainji}
\item Blench's review of East Kainji languages (a lead for Vori-specific material; I did not read it): \url{https://www.academia.edu/38612586}
\end{itemize}

\section{Religious and Community Resources}
\begin{itemize}
\item Joshua Project prayer card for the Vori (population 3,900): \url{https://joshuaproject.net/pray/cards/single/15138/NI}
\end{itemize}
No scripture translation was confirmed. Global Recordings lists Voro (a different language), so it is not cited here.

\section{Digital and Audio Resources}
No Vori-specific audio, video or corpus was found. The ELDP Kainji documentation collection (\url{https://eldp.access.preservica.com/dk0145/}) did not appear to include Vori.

\section{Curriculum Development Context}
Vori is essentially undocumented for educational purposes. A curriculum would need a documentation phase first: a wordlist, an orthography agreed with the community, and recorded texts. Blench's work on Kainji languages is the likeliest starting point for contacts and data.
\end{document}